\documentclass[aspectratio=169]{beamer}
\usetheme{Madrid}
\usecolortheme{default}
\usepackage[utf8]{inputenc}
\usepackage{amsmath}
\usepackage{amsfonts}
\usepackage{amssymb}

\setbeamertemplate{navigation symbols}{}
\setbeamerfont{block body}{size=\small}
\setbeamerfont{alertblock body}{size=\small}

\title{CRISC Domain B -- 100 Exam-Style Questions (Part 3)}
\subtitle{IT Risk Assessment -- Fresh Scenarios (Answer on the Next Slide)}
\author{CRISC Mastery Series}
\date{\today}

% ------------------------------------------------------------
% \hq{number}{topic}{question}{options}{letter}{explanation}{trap}
% ------------------------------------------------------------
\newcommand{\hq}[7]{%
	\begin{frame}{Q#1: #2}
		\begin{block}{Question}
			#3
		\end{block}
		\begin{enumerate}[(A)]
			#4
		\end{enumerate}
	\end{frame}
	\begin{frame}{Q#1: Explanation}
		\begin{block}{Correct Answer: #5}
			\begin{itemize}
				#6
			\end{itemize}
		\end{block}
		\begin{alertblock}{Critical Trap}
			#7
		\end{alertblock}
	\end{frame}
}

\begin{document}
	
	\begin{frame}\titlepage\end{frame}
	
	\begin{frame}{Instructions}
		\begin{alertblock}{How to Use This Deck}
			Think like a risk practitioner: business objectives first, owners decide, assess before acting.\\
			Read for BEST, FIRST and MOST. Answer on the slide after each question.
		\end{alertblock}
	\end{frame}
	
	\begin{frame}{Outline}
		\tableofcontents
	\end{frame}
	
	% ============================================================
	\section{Risk Identification}
	% ============================================================
	
	\hq{1}{Regulatory Horizon}
	{A regulator publishes a consultation paper on AI risk that may become law in 18 months. What should the risk practitioner do FIRST?}
	{\item Wait until the law is passed
		\item Ignore it as non-binding
		\item Analyze potential exposure and monitor developments
		\item Rewrite all AI policies today}
	{C}
	{\item Consultation signals future obligations
		\item Early analysis enables planning
		\item Rewriting policies before certainty wastes effort
		\item Ignoring the horizon creates late surprises}
	{Consultations are early warnings, not certainties.}
	
	\hq{2}{Board Risk Appetite}
	{The board states it has a LOW appetite for data breaches but a MODERATE appetite for service interruption. What does this mean for the risk assessment scope?}
	{\item Treat all risks equally
		\item Prioritize breach-related scenarios and use lighter scrutiny on brief outages
		\item Ignore service interruption entirely
		\item Focus only on fire risks}
	{B}
	{\item Appetite differentials shape priorities
		\item Breach scenarios warrant deeper analysis
		\item Service risks still assessed but proportionate to appetite
		\item Scope should reflect the appetite profile}
	{Scope follows appetite differentials, not uniformity.}
	
	\hq{3}{Business Model Change}
	{A firm shifts from one-time product sales to a subscription model. Which new risk area should be identified FIRST?}
	{\item Office supply risk
		\item Customer data lifecycle and recurring billing availability
		\item Printer maintenance risk
		\item Parking risk}
	{B}
	{\item Subscription models depend on recurring billing and customer data
		\item Data lifecycle drives privacy and retention obligations
		\item Billing availability directly impacts revenue
		\item Physical risks remain but are secondary to the model}
	{New business models create new risk categories.}
	
	\hq{4}{Employee Turnover}
	{High turnover in a security operations team creates which category of risk?}
	{\item Equipment depreciation risk
		\item Loss of institutional knowledge and control effectiveness degradation
		\item Increase in advertising cost
		\item Reduced tax liability}
	{B}
	{\item Knowledge walks out with people
		\item Controls depend on skilled operators
		\item Replacement time extends exposure}
	{People risk is a control risk.}
	
	\hq{5}{Post-Merger Integration}
	{After a merger, the acquired company uses unsupported software. What should the practitioner do FIRST?}
	{\item Ignore legacy systems
		\item Assess criticality of those systems and identify compensating controls
		\item Migrate immediately without planning
		\item File a legal complaint}
	{B}
	{\item Unsupported software creates vulnerabilities
		\item Criticality defines urgency
		\item Compensating controls bridge the gap
		\item Immediate migration without planning increases risk}
	{Assess and compensate before mass migration.}
	
	\hq{6}{Environmental Scanning}
	{The risk practitioner subscribes to threat intelligence but never uses it. What is the BEST improvement?}
	{\item Cancel the subscription
		\item Integrate intelligence into a routine risk identification process
		\item Increase the subscription volume
		\item Delegate reading to staff without a process}
	{B}
	{\item Intelligence must feed a process
		\item Routine integration prevents forgotten alerts
		\item Volume without process adds noise}
	{Intelligence unused is intelligence wasted.}
	
	\hq{7}{External Dependency}
	{The organization relies on a single public DNS provider. What category of risk should be identified?}
	{\item Property tax risk
		\item External dependency and concentration risk
		\item Employee parking risk
		\item Supplier branding risk}
	{B}
	{\item Single external dependency means availability risk
		\item Concentration amplifies impact
		\item DNS is critical to customer-facing services}
	{External dependencies create availability risk.}
	
	\hq{8}{Fraud Scheme}
	{Finance staff can create and approve vendor payments. What risk category is MOST relevant?}
	{\item Fraud and segregation-of-duties risk
		\item Advertising cost risk
		\item Software licensing risk
		\item Training cost risk}
	{A}
	{\item Combined create/approve authority creates fraud exposure
		\item SoD prevents this
		\item Financial controls address the gap}
	{Combined duties enable fraud.}
	
	\hq{9}{New Location}
	{Opening an office in a new country raises which FIRST-order risk identification concern?}
	{\item Local labor laws, data residency, and regulatory differences
		\item Office coffee brand
		\item Local advertisement style
		\item Parking spot count}
	{A}
	{\item Legal and regulatory differences are primary
		\item Data residency affects architecture
		\item Labor law affects HR and data}
	{Geography brings legal and data risks.}
	
	\hq{10}{Major Outage}
	{A cloud provider experiences a regional outage affecting several customers. What should the practitioner do NEXT?}
	{\item Switch to another provider immediately
		\item Assess whether the organization has similar exposure and revisit single-provider assumptions
		\item Ignore it
		\item Cancel the contract}
	{B}
	{\item External events reveal your own exposure
		\item Revisit concentration assumptions
		\item Immediate switching may be infeasible or unnecessary}
	{External events reveal internal exposure.}
	
	\hq{11}{Vendor Change}
	{A critical vendor changes its data-processing location without notice. What risk is most relevant?}
	{\item Data residency and legal jurisdiction
		\item Office rent
		\item Travel expense
		\item Holiday schedule}
	{A}
	{\item Data location affects compliance and legal exposure
		\item Jurisdictional issues can trigger regulatory risk
		\item Location changes require re-assessment}
	{Data location is a regulatory and legal concern.}
	
	\hq{12}{AI Model}
	{A new AI model is used to make credit decisions. What risk should be identified FIRST?}
	{\item Model bias, explainability, and regulatory compliance
		\item Office supply risk
		\item Printer risk
		\item Coffee risk}
	{A}
	{\item Bias creates discrimination risk
		\item Explainability affects regulatory compliance
		\item Model failures cascade into credit decisions}
	{AI risks include bias, explainability, and compliance.}
	
	\hq{13}{Weak Password Culture}
	{Employees reuse the same password across work and personal accounts. What risk category should be identified?}
	{\item Credential compromise and lateral movement risk
		\item Air conditioning risk
		\item Fire suppression risk
		\item Building security risk}
	{A}
	{\item Reuse spreads compromise risk
		\item Lateral movement follows credential theft
		\item Compensating controls include MFA and credential monitoring}
	{Credential reuse enables widespread compromise.}
	
	\hq{14}{Automation}
	{A process is automated without human review. What risk should be identified?}
	{\item Silent errors scaling automatically
		\item Travel risk
		\item Catering risk
		\item Office layout risk}
	{A}
	{\item Automated errors scale
		\item Human review catches mistakes
		\item Monitoring and exception reporting are key controls}
	{Automation scales both efficiency and errors.}
	
	\hq{15}{Data Quality}
	{Reports depend on inconsistent data sources. What risk should be identified?}
	{\item Decision risk from inaccurate data
		\item Fire risk
		\item Library risk
		\item Paving risk}
	{A}
	{\item Inaccurate data drives wrong decisions
		\item Data quality is a business risk
		\item Lineage and validation reduce it}
	{Bad data creates bad decisions.}
	
	% ============================================================
	\section{Risk Analysis}
	% ============================================================
	
	\hq{16}{Correlation Assumption}
	{Analysts assume two risks are independent. Evidence suggests they co-occur. What is the impact on the analysis?}
	{\item No impact
		\item Combined risk may be significantly understated
		\item Risk decreases
		\item ALE becomes zero}
	{B}
	{\item Independence underestimates joint probability
		\item Co-occurrence amplifies exposure
		\item Recalculate using correlation}
	{Independence assumptions understate correlated risk.}
	
	\hq{17}{Qualitative Scale}
	{The organization uses High/Medium/Low but has no definitions. What is the BEST improvement?}
	{\item Remove the scale
		\item Define each level with measurable criteria
		\item Add more levels without definitions
		\item Use only colors}
	{B}
	{\item Definitions reduce subjectivity
		\item Measurable criteria improve consistency
		\item More levels without meaning add noise}
	{Define what each level means.}
	
	\hq{18}{Expert Calibration}
	{Experts consistently underestimate likelihood. What is the BEST corrective measure?}
	{\item Ask experts to estimate faster
		\item Calibrate estimates with historical data and independent challenge
		\item Average expert estimates
		\item Eliminate expert judgment}
	{B}
	{\item Data calibrates optimism bias
		\item Independent challenge reduces anchoring
		\item Fast estimates increase error}
	{Data and challenge calibrate expert estimates.}
	
	\hq{19}{Data Sampling}
	{A large transaction dataset is analyzed. Which approach is MOST defensible for extrapolation?}
	{\item Selecting transactions alphabetically
		\item A statistically valid sample with documented methodology
		\item Choosing transactions that look interesting
		\item Sampling only recent transactions}
	{B}
	{\item Statistical sampling supports defensible extrapolation
		\item Documented methodology improves auditability
		\item Biased selection invalidates conclusions}
	{Sampling must be statistically defensible.}
	
	\hq{20}{New System}
	{A new system has no historical data. What is the BEST analysis technique?}
	{\item Wait two years for data
		\item Scenario-based analysis with expert judgment and benchmarks
		\item Use last year\textquotesingle s numbers from another system
		\item Skip analysis}
	{B}
	{\item Scenarios and benchmarks suit new systems
		\item Expert judgment supplements missing data
		\item Borrowing unrelated data misleads}
	{New systems need scenarios, not fake data.}
	
	\hq{21}{Model Overfit}
	{A model fits past incidents perfectly but fails on new cases. What is the issue?}
	{\item Overfitting
		\item Underfitting
		\item Determinism
		\item Randomness}
	{A}
	{\item Overfitting memorizes past, not future
		\item Simpler, validated models generalize better
		\item Cross-validation reveals overfit}
	{Perfect fit on old data often fails on new data.}
	
	\hq{22}{ALE with Downtime}
	{An hour of downtime costs \$50K. A risk has 20\% annual likelihood of a 4-hour outage. What is the ALE?}
	{\item \$40K
		\item \$10K
		\item \$100K
		\item \$200K}
	{A}
	{\item Impact $= 4 \times 50K = \$200K$
		\item Likelihood $= 20\%$
		\item ALE $= 200K \times 0.2 = \$40K$
		\item \$200K is single-loss impact}
	{ALE equals impact multiplied by likelihood.}
	
	\hq{23}{Sensitive Assumption}
	{The rating changes dramatically when the assumed detection time doubles. What should the practitioner do?}
	{\item Ignore the sensitivity
		\item Investigate detection capability and refine the analysis
		\item Increase the rating to High
		\item Report the highest possible figure}
	{B}
	{\item Sensitive assumptions deserve validation
		\item Detection data improves accuracy
		\item Arbitrary changes mislead}
	{Test and refine the sensitive assumption.}
	
	\hq{24}{Baseline Drift}
	{The risk profile gradually changes without any trigger event. What monitoring helps detect drift?}
	{\item Annual snapshot only
		\item Trend analysis over time
		\item None
		\item Marketing surveys}
	{B}
	{\item Trends reveal gradual change
		\item Snapshots miss movement between them
		\item Continuous or periodic trend analysis catches drift}
	{Drift is invisible to single snapshots.}
	
	\hq{25}{Worst-Case Use}
	{When is worst-case analysis MOST useful?}
	{\item Routine operations
		\item Business continuity and insurance decisions
		\item Marketing budget decisions
		\item Staff scheduling}
	{B}
	{\item Worst-case informs continuity and insurance
		\item Routine decisions use expected values
		\item Marketing and staff use different measures}
	{Worst case informs resilience planning.}
	
	\hq{26}{Data Fusion}
	{Internal incident data and threat intelligence conflict. What is the BEST approach?}
	{\item Use only internal data
		\item Reconcile by considering both sources and validate assumptions
		\item Use only intelligence
		\item Average the two}
	{B}
	{\item Reconciliation improves accuracy
		\item Both sources have value
		\item Averages hide meaning}
	{Reconcile, do not average.}
	
	\hq{27}{Data Silos}
	{Each department keeps its own risk spreadsheet. What is the MAIN analytical risk?}
	{\item Increased visibility
		\item Duplicated and inconsistent risk data
		\item Better ownership
		\item Faster reporting}
	{B}
	{\item Silos duplicate and conflict
		\item Consistency requires a central register
		\item Visibility increases with centralization}
	{Siloed spreadsheets create inconsistency.}
	
	\hq{28}{Analysis Purpose}
	{Risk analysis is BEST used to:}
	{\item Eliminate the need for judgment
		\item Provide structured information for decisions
		\item Replace governance
		\item Avoid escalation}
	{B}
	{\item Analysis informs decisions
		\item Judgment and governance remain essential
		\item Escalation remains part of the process}
	{Analysis informs; decisions remain with people.}
	
	\hq{29}{Change Threshold}
	{The business changes its impact definitions after the assessment. What should happen?}
	{\item Nothing
		\item Re-score affected risks using the new definitions
		\item Ignore the new definitions
		\item Delete the assessment}
	{B}
	{\item Definitions drive scoring
		\item Consistent definitions keep ratings meaningful
		\item Ignoring definitions invalidates results}
	{When definitions change, ratings must follow.}
	
	\hq{30}{Residual Estimation}
	{Residual risk estimation requires:}
	{\item Inherent risk plus budget
		\item Control effectiveness evidence
		\item Marketing insight
		\item Vendor invoices}
	{B}
	{\item Residual depends on proven control effectiveness
		\item Evidence avoids guesswork
		\item Budget and invoices do not estimate risk}
	{Residual requires control evidence.}
	
	% ============================================================
	\section{Risk Evaluation and Appetite}
	% ============================================================
	
	\hq{31}{Tiered Appetite}
	{The organization sets a different appetite for customer data than internal data. What should the practitioner do?}
	{\item Apply the strictest appetite everywhere
		\item Reflect the tiered appetite in risk ratings and thresholds
		\item Ignore the difference
		\item Use only the weakest appetite}
	{B}
	{\item Tiering matches sensitivity
		\item Thresholds should reflect the tier
		\item Applying strictest everywhere wastes effort}
	{Appetite tiers drive thresholds.}
	
	\hq{32}{Appetite Escalation}
	{Business expansion inherently increases risk above appetite. What is the BEST action?}
	{\item Hide the increase
		\item Escalate for formal review and appetite reassessment
		\item Approve without escalation
		\item Cancel the expansion}
	{B}
	{\item Higher risk above appetite requires review
		\item Leadership may adjust appetite
		\item Silent proceeding erodes governance}
	{Escalate when appetite conflicts with strategy.}
	
	\hq{33}{Risk Tolerance in Incident}
	{An incident exceeded tolerance. What should happen FIRST?}
	{\item Close the incident without further action
		\item Investigate root cause and reassess exposure
		\item Change tolerance thresholds
		\item Increase the budget}
	{B}
	{\item Root cause informs needed action
		\item Tolerance breaches signal governance failures
		\item Raising thresholds hides the signal}
	{Breach triggers investigation, not threshold shift.}
	
	\hq{34}{Residual Escalation}
	{Residual risk after treatment is above appetite. Who decides next?}
	{\item The controls team
		\item The vendor
		\item The authorized authority per governance
		\item The auditor}
	{C}
	{\item Above-appetite decisions require authority
		\item Options are more treatment or acceptance
		\item Vendor and audit do not own business risk}
	{Above appetite goes to the appropriate authority.}
	
	\hq{35}{Acceptance Duration}
	{A formal risk acceptance should usually include which element?}
	{\item A hidden expiration
		\item An open-ended acceptance with no review
		\item A review date or expiration to trigger reassessment
		\item Only the analyst\textquotesingle s signature}
	{C}
	{\item Review dates prevent stale acceptance
		\item Open-ended acceptance hides drift
		\item Owners must re-commit}
	{Acceptance must expire and be revisited.}
	
	\hq{36}{Appetite Communication}
	{Employees do not know the risk appetite. What is the BEST improvement?}
	{\item Post the appetite once per year
		\item Embed appetite thresholds into operating procedures and decisions
		\item Add more forms
		\item Ignore it}
	{B}
	{\item Embedded thresholds guide daily choices
		\item Documentation alone is forgotten
		\item Procedures make appetite practical}
	{Embed appetite in operations.}
	
	\hq{37}{Capacity and Strategy}
	{Strategic growth may breach capacity. What is the BEST action?}
	{\item Ignore capacity
		\item Model capacity impact and plan treatment before growth
		\item Freeze all growth
		\item Accept risk without analysis}
	{B}
	{\item Capacity breaches threaten existence
		\item Proactive treatment supports growth
		\item Freezing everything avoids opportunity}
	{Plan before capacity is breached.}
	
	\hq{38}{Loss of Appetite Confidence}
	{A rating was revised after political pressure from a business unit. What is the BEST response?}
	{\item Accept the revision
		\item Challenge pressure, document the independent basis and escalate if needed
		\item Increase all ratings
		\item Ignore the issue}
	{B}
	{\item Independence protects integrity
		\item Documentation defends decisions
		\item Escalation addresses pressure}
	{Pressure must not drive ratings.}
	
	\hq{39}{Appetite Drift}
	{Appetite was set five years ago and never revisited. What is the most likely issue?}
	{\item No issue
		\item Appetite may no longer reflect the business context
		\item Appetite is too high
		\item Appetite is too low}
	{B}
	{\item Business context changes
		\item Appetite must be refreshed periodically
		\item Stale appetite misguides decisions}
	{Appetite must be revisited with strategy.}
	
	\hq{40}{Post-Acceptance Monitoring}
	{An accepted risk reappears in an audit finding. What should the practitioner do?}
	{\item Ignore the finding
		\item Verify whether the acceptance is documented and whether conditions changed
		\item Withdraw the finding
		\item Blame the auditor}
	{B}
	{\item Documented acceptance defends decisions
		\item Changed conditions may invalidate the original acceptance
		\item Audit findings should be reviewed, not withdrawn}
	{Verify the acceptance and current conditions.}
	
	% ============================================================
	\section{Risk Response}
	% ============================================================
	
	\hq{41}{Transfer Consequences}
	{Insurance covers financial loss but not downtime. What does this mean for the residual risk?}
	{\item All risk is eliminated
		\item Financial residual is covered but operational residual remains
		\item Residual is fully transferred
		\item Residual is zero}
	{B}
	{\item Transfer is financial only
		\item Operational and reputational risk remain
		\item Plan controls for operational impact}
	{Transfer is financial, not operational.}
	
	\hq{42}{Cheaper Alternative}
	{The first control exceeds the expected loss. What should the practitioner explore FIRST?}
	{\item Implement it anyway
		\item Search for less expensive controls with similar effect
		\item Accept the risk without review
		\item Cancel the project}
	{B}
	{\item Alternatives may be more efficient
		\item Cost-benefit favors proportionate controls
		\item Unilateral acceptance without review is premature}
	{Explore cheaper options before accepting.}
	
	\hq{43}{Data Custody}
	{Data stored with a vendor is encrypted by the organization. Who is accountable for the data?}
	{\item The vendor
		\item The organization
		\item The regulator
		\item The insurance provider}
	{B}
	{\item Accountability remains with the data owner
		\item Encryption supports but does not shift accountability
		\item Regulators and insurers do not take accountability}
	{Accountability follows the data owner.}
	
	\hq{44}{Response Tradeoff}
	{A response reduces likelihood but increases impact if it fails. What should the practitioner evaluate?}
	{\item Ignore the new impact
		\item The net effect on residual risk and new risks introduced
		\item Only the likelihood reduction
		\item Only implementation cost}
	{B}
	{\item Net residual and new risk matter
		\item Secondary risks must be assessed
		\item Partial views mislead decisions}
	{Assess the net effect and secondary risks.}
	
	\hq{45}{Avoidance Cost}
	{Terminating a service avoids its risk. What else must be considered?}
	{\item Only the risk removed
		\item Opportunity cost, contractual obligations and stakeholder impact
		\item Only the cost of termination
		\item Only legal fees}
	{B}
	{\item Avoidance has opportunity and stakeholder costs
		\item Contracts may prevent immediate exit
		\item Full impact analysis is required}
	{Avoidance is not free; consider lost value.}
	
	\hq{46}{Preventive Gap}
	{An attack succeeded despite preventive controls. What does this suggest?}
	{\item Prevention is unnecessary
		\item Prevention should be layered with detective and corrective controls
		\item The risk is now zero
		\item The controls were useless}
	{B}
	{\item Layered controls support each other
		\item Prevention will not stop everything
		\item Detection and correction reduce impact}
	{Prevention needs detection and correction layers.}
	
	\hq{47}{Detection Latency}
	{Detection latency is too long. What is the BEST improvement?}
	{\item Add more prevention only
		\item Tune detection, improve ingestion and alerting
		\item Increase the budget without analysis
		\item Cancel detection}
	{B}
	{\item Detection tuning reduces latency
		\item Better ingestion improves coverage
		\item Prevention and detection complement}
	{Improve detection, not just prevention.}
	
	\hq{48}{Compensating Adequacy}
	{A compensating control reduces risk but not below the original target. What should happen?}
	{\item Accept it
		\item Escalate because it does not achieve the risk objective
		\item Ignore the gap
		\item Remove the primary control}
	{B}
	{\item Compensation must meet the risk target
		\item Escalation is required if not achieved
		\item Acceptance needs authorized decision}
	{Compensation must actually close the gap.}
	
	\hq{49}{Response Ownership}
	{Who is accountable for ensuring a risk response is implemented?}
	{\item The project sponsor
		\item The risk owner
		\item The auditor
		\item The helpdesk}
	{B}
	{\item Risk owners are accountable for responses
		\item Sponsors and auditors support but do not own the risk
		\item Helpdesk does not own business risk}
	{Risk owner oversees response execution.}
	
	\hq{50}{Verification After Change}
	{A major infrastructure change occurred. What should the practitioner do NEXT?}
	{\item Assume risk unchanged
		\item Reassess risk and verify control effectiveness after change
		\item Skip verification
		\item Rely on old reports}
	{B}
	{\item Changes alter risk
		\item Post-change verification confirms controls operate
		\item Old reports describe a different environment}
	{Re-verify after major changes.}
	
	\hq{51}{Formal vs Informal Acceptance}
	{A risk is informally tolerated by staff. What is the BEST action?}
	{\item Accept it as-is
		\item Formalize the acceptance with the risk owner or treat the risk
		\item Ignore it
		\item Increase the budget}
	{B}
	{\item Informal acceptance has no accountability
		\item Formal acceptance creates governance
		\item Treatment reduces exposure}
	{Informal tolerance is not acceptance.}
	
	\hq{52}{Assurance Alternatives}
	{A vendor provides no third-party assurance. What should the practitioner do?}
	{\item Rely on marketing claims
		\item Perform direct assessment or require compensating controls
		\item Skip the vendor
		\item Sign without review}
	{B}
	{\item Direct assessment provides assurance
		\item Compensating controls reduce exposure
		\item Marketing claims are not assurance}
	{Assurance requires evidence.}
	
	\hq{53}{Exception Exit}
	{A policy exception is no longer required. What should be done?}
	{\item Keep it open
		\item Close it and update records
		\item Renew it automatically
		\item Ignore it}
	{B}
	{\item Closing reduces unnecessary exposure
		\item Records should reflect the current state
		\item Automatic renewal keeps exceptions alive}
	{Close exceptions that are no longer needed.}
	
	\hq{54}{Response Chain}
	{A control depends on a vendor that itself depends on a subcontractor. What is the KEY risk?}
	{\item Fourth-party dependency and loss of visibility
		\item Marketing cost
		\item Office lease
		\item Catering}
	{A}
	{\item Fourth-party risk affects control reliability
		\item Visibility is reduced through layers
		\item Chain failures disrupt operations}
	{Chain dependencies require extended oversight.}
	
	\hq{55}{Response Justification}
	{Which best justifies a mitigation investment?}
	{\item Risk is low
		\item Risk is above appetite and controls are cost-effective
		\item Risk has not occurred
		\item Risk is very unlikely}
	{B}
	{\item Above-appetite risk requires action
		\item Cost-effectiveness protects value
		\item Past absence is not current risk}
	{Mitigate when above appetite and cost-effective.}
	
	% ============================================================
	\section{Controls}
	% ============================================================
	
	\hq{56}{Design Document}
	{A control is described only in a slide with no procedure. What is the deficiency type?}
	{\item Design deficiency
		\item Operating deficiency
		\item Independence issue
		\item Scope issue}
	{A}
	{\item Without a procedure, the control lacks design
		\item Design must include roles and steps
		\item Operating issues assume a proper design}
	{Design requires procedures, not slides.}
	
	\hq{57}{Defense Layers}
	{Which combination reduces the risk of common-mode failure?}
	{\item Duplicate firewalls from the same vendor
		\item Diverse network isolation, host controls, and monitoring
		\item Only perimeter controls
		\item Only endpoint controls}
	{B}
	{\item Diversity limits common-mode failure
		\item Perimeter and endpoint only leave gaps
		\item Duplicates can share the same flaw}
	{Diversity limits common-mode failure.}
	
	\hq{58}{Owner Evidence}
	{Which evidence BEST supports a control operating effectively?}
	{\item A policy exists
		\item The control was executed with evidence and exceptions reviewed
		\item A budget was allocated
		\item A team was formed}
	{B}
	{\item Execution evidence proves operation
		\item Policies and budgets are inputs
		\item Team formation does not prove operation}
	{Execution evidence proves effectiveness.}
	
	\hq{59}{Correction Effectiveness}
	{A corrective control restores service but loses 12 hours of data with an RPO of 4 hours. What does this mean?}
	{\item The corrective control is adequate
		\item The corrective control does not meet the RPO
		\item The RPO should be raised
		\item No issue exists}
	{B}
	{\item RPO violation indicates data loss risk
		\item Corrective controls must support RPO
		\item RPO changes require business approval}
	{Correction must respect recovery objectives.}
	
	\hq{60}{Audit Independence}
	{An internal audit report on IT controls was written by the person who built the controls. What is the MAIN issue?}
	{\item Independence is compromised
		\item The report is too long
		\item The report costs too much
		\item The report lacks visuals}
	{A}
	{\item Independence affects objectivity
		\item Self-review produces bias
		\item Independent review is needed}
	{Self-assessment is not independent audit.}
	
	\hq{61}{Gap Prioritization}
	{Multiple gaps are found. What MOST drives prioritization?}
	{\item The number of findings
		\item Exposure, likelihood and business impact
		\item The cheapest fixes
		\item The oldest gaps}
	{B}
	{\item Exposure drives priority
		\item Impact aligns with business objectives
		\item Age and cost are secondary}
	{Prioritize by exposure, not count.}
	
	\hq{62}{Automation Monitoring}
	{An automated control processes payments. What monitoring is essential?}
	{\item None, it is automated
		\item Exception monitoring, logging and periodic validation
		\item Only annual review
		\item Only manual override}
	{B}
	{\item Automation needs monitoring
		\item Exceptions reveal failure
		\item Periodic validation confirms operation}
	{Automate with visibility.}
	
	\hq{63}{Proportional Controls}
	{Which principle guides control proportionality?}
	{\item All controls must be equal
		\item Controls should match the risk level and appetite
		\item Controls must be automatic
		\item Controls must be manual}
	{B}
	{\item Proportionality aligns effort with risk
		\item Low-risk controls should not exceed value
		\item High-risk controls need stronger measures}
	{Controls should fit the risk size.}
	
	\hq{64}{Redundant Justified}
	{When is redundancy MOST justified?}
	{\item When the risk is low
		\item When common-mode failure could cause severe impact
		\item When the budget allows
		\item When it looks impressive}
	{B}
	{\item Severe impact justifies redundancy
		\item Low risk rarely justifies duplicated cost
		\item Budget is not the driver}
	{Redundancy protects against severe impact.}
	
	\hq{65}{CSA Validation}
	{What adds credibility to a control self-assessment?}
	{\item Independence and evidence sampling
		\item More self-reports
		\item Faster completion
		\item No follow-up}
	{A}
	{\item Independent validation reduces bias
		\item Evidence sampling supports conclusions
		\item Rushed CSAs lack rigor}
	{CSA credibility requires validation.}
	
	% ============================================================
	\section{Monitoring and Reporting}
	% ============================================================
	
	\hq{66}{Lagging Indicator}
	{Which is a LAGGING indicator of phishing exposure?}
	{\item Number of employees who failed last quarter\textquotesingle s phishing simulation
		\item Click rate trends
		\item Email gateway rules
		\item Training attendance}
	{A}
	{\item Past simulation failures are lagging
		\item Trends and configurations are inputs
		\item Attendance is an input}
	{Lagging indicators record what already happened.}
	
	\hq{67}{Alert Fatigue}
	{Analysts ignore most alerts. What is the BEST fix?}
	{\item Increase alert volume
		\item Reduce and prioritize alerts with clear response playbooks
		\item Cancel monitoring
		\item Delegate to the helpdesk}
	{B}
	{\item Prioritization improves response quality
		\item Playbooks standardize action
		\item Cancelling removes visibility}
	{Fewer, better alerts beat more, ignored alerts.}
	
	\hq{68}{Regulator Reporting}
	{Which report format BEST supports a regulatory request?}
	{\item A marketing brochure
		\item A short executive dash
		\item A structured evidence-based report tied to requirements
		\item An email summary}
	{C}
	{\item Regulators need evidence and traceability
		\item Structured reporting supports auditability
		\item Marketing and email are inadequate}
	{Regulators need evidence and structure.}
	
	\hq{69}{Register Changes}
	{Which change should trigger an immediate register update?}
	{\item A new critical vendor was onboarded
		\item A staff meeting was scheduled
		\item A brochure was printed
		\item A holiday was announced}
	{A}
	{\item Vendor onboarding changes exposure
		\item Critical vendors introduce third-party risk
		\item Administrative events do not change risk}
	{Onboarding critical vendors changes exposure.}
	
	\hq{70}{Real-Time Limits}
	{Which scenario does NOT justify real-time monitoring?}
	{\item High-value transaction integrity
		\item Detection of network intrusions
		\item Annual marketing plan progress
		\item Malware outbreak detection}
	{C}
	{\item Marketing plan progress is periodic
		\item Transaction, intrusion, and malware risks need real-time
		\item Periodic reviews suffice for strategic plans}
	{Real-time is for urgent risks.}
	
	\hq{71}{KRI Overload}
	{The dashboard has 200 KRIs. What is the BEST remedy?}
	{\item Add more visuals
		\item Reduce to a focused set tied to appetite with clear ownership
		\item Remove all KRIs
		\item Increase the budget}
	{B}
	{\item Focused metrics drive decisions
		\item Appetite alignment gives meaning
		\item Removing KRIs loses signal}
	{Fewer, owned KRIs improve response.}
	
	\hq{72}{Indicator Sensitivity}
	{KRIs rarely change even when risk changes. What is the issue?}
	{\item Insensitive thresholds
		\item Excellent monitoring
		\item Correct configuration
		\item Data volume}
	{A}
	{\item Insensitivity hides risk change
		\item Thresholds must track meaningful movement
		\item Stable indicators may be miscalibrated}
	{Insensitive KRIs hide risk.}
	
	\hq{73}{Green Is Not Good}
	{The board interprets green as no risk. What should the practitioner explain?}
	{\item Green means risk is within appetite, not absent
		\item Green means risk is zero
		\item Green means no action needed
		\item Green means budget is sufficient}
	{A}
	{\item Green means within tolerance
		\item Risk still exists
		\item Monitoring continues}
	{Green means within appetite, not risk-free.}
	
	\hq{74}{Messaging Style}
	{Which phrase BEST supports executive risk discussion?}
	{\item We patched 400 servers
		\item Our breach exposure this quarter decreased 15\% and remains within appetite
		\item The SIEM correlated 500 events
		\item The firewall vendor called}
	{B}
	{\item Business language drives decisions
		\item Metrics support trend discussion
		\item Technical and anecdotal statements distract}
	{Business language drives decisions.}
	
	\hq{75}{Root Cause of Miss}
	{An incident went undetected for 48 hours. What should be reviewed FIRST?}
	{\item Marketing plans
		\item Detection coverage, logging and alerting
		\item Vendor invoices
		\item Budget spreadsheets}
	{B}
	{\item Detection latency points to monitoring gaps
		\item Coverage and alerting are direct controls
		\item Other items are not related}
	{Undetected incidents suggest monitoring gaps.}
	
	\hq{76}{Register Quality}
	{Register entries lack owners and dates. What is the MAIN impact?}
	{\item Ownership and tracking weaken, reducing accountability
		\item Reporting improves
		\item Costs decrease
		\item Controls operate better}
	{A}
	{\item Owners and dates are core accountability fields
		\item Missing them creates governance gaps
		\item Tracking requires structure}
	{Core fields enable accountability.}
	
	\hq{77}{Threshold Ownership}
	{Escalation thresholds exist but no one knows who to notify. What should be added?}
	{\item Named recipients and contact paths per scenario
		\item A weekly newsletter
		\item More thresholds
		\item Another tool}
	{A}
	{\item Escalation requires recipients
		\item Clarity improves response time
		\item Tools alone do not provide ownership}
	{Escalation needs recipients, not more rules.}
	
	\hq{78}{Data Provenance}
	{KRI data has unclear sources. What is the BEST fix?}
	{\item Document data sources and validate integrity
		\item Use the data anyway
		\item Average across sources
		\item Delete the KRIs}
	{A}
	{\item Provenance supports trust
		\item Validation confirms accuracy
		\item Unclear sources mislead}
	{Provenance builds KRI trust.}
	
	\hq{79}{Audit Findings}
	{Audit findings are repeated year over year. What does this suggest?}
	{\item Findings are strong
		\item Remediation is ineffective
		\item Findings are unimportant
		\item The auditor is wrong}
	{B}
	{\item Repeat findings signal unmet remediation
		\item Ownership and accountability are weak
		\item Escalation may be required}
	{Repeated findings signal failed remediation.}
	
	\hq{80}{Directional Risk}
	{A KRI improves but stress indicators elsewhere worsen. What is the BEST view?}
	{\item Ignore stress indicators
		\item Combine KRI direction with other signals for a full picture
		\item Use only the best KRI
		\item Cancel the dashboard}
	{B}
	{\item Single indicators mislead
		\item Combined signals reveal net direction
		\item Selective metrics distort decisions}
	{Combine signals for a true direction.}
	
	% ============================================================
	\section{BIA, Recovery and Third Parties}
	% ============================================================
	
	\hq{81}{BIA Scope}
	{A BIA covers the revenue-critical order system but omits the payroll system. What is the MAIN risk?}
	{\item Payroll risks are unassessed and may cause hidden business impact
		\item Revenue systems are irrelevant
		\item No issue
		\item HR is not relevant}
	{A}
	{\item Impact extends to employee trust and compliance
		\item BIA must include all critical processes
		\item Omitting payroll hides impact}
	{BIA must reflect all critical processes.}
	
	\hq{82}{RTO Cascade}
	{A critical app has a 2-hour RTO. It depends on a database with a 6-hour RTO. What is the issue?}
	{\item No issue
		\item The dependency prevents meeting the RTO
		\item The database RTO is too short
		\item The app RTO is too long}
	{B}
	{\item Dependent recovery times must align
		\item The database recovery time blocks app recovery
		\item Misalignment is a common BIA failure}
	{Dependencies must support the primary RTO.}
	
	\hq{83}{RPO Cascade}
	{Data flows through three systems with RPOs of 1, 4, and 8 hours. What is the effective RPO?}
	{\item 1 hour
		\item 13 hours
		\item 8 hours (the weakest link)
		\item 4 hours}
	{C}
	{\item The weakest link governs data loss
		\item Each system can lose up to its RPO
		\item Effective RPO is the worst RPO in the chain}
	{The weakest link determines effective RPO.}
	
	\hq{84}{MTD Type}
	{Which definition fits an MTD of 12 hours?}
	{\item No more than 12 hours of data loss
		\item Total outage beyond 12 hours causes unacceptable impact
		\item Restore within 12 hours
		\item 12 hours is the average recovery time}
	{B}
	{\item MTD is the impact threshold
		\item RTO is the recovery target
		\item RPO is about data loss}
	{MTD is the impact threshold.}
	
	\hq{85}{BIA Ownership}
	{The BIA was written by IT without business input. What is the MAIN concern?}
	{\item Business impact may be misjudged
		\item IT budgets may be cut
		\item The document is too long
		\item The process is too fast}
	{A}
	{\item Business owners know impact
		\item IT understands capability, not business value
		\item Misjudged impact skews recovery}
	{BIA needs business owners.}
	
	\hq{86}{Backup Integrity}
	{Backups exist but are never tested. What is the KEY risk?}
	{\item Backups may fail on restore
		\item The cost of backups
		\item Backup speed
		\item Backup media branding}
	{A}
	{\item Untested backups may not restore
		\item Restore testing is essential
		\item Success relies on verification}
	{Untested backups are unreliable backups.}
	
	\hq{87}{DR Test Failure}
	{A DR test fails on a dependency not covered by the plan. What is the BEST action?}
	{\item Ignore the failure
		\item Update the DR plan to include the dependency and retest
		\item Raise the RTO
		\item Cancel tests}
	{B}
	{\item Dependency gaps are common failures
		\item Plan updates close gaps
		\item RTO changes require business agreement}
	{Fix the plan, retest the gap.}
	
	\hq{88}{Vendor Assessment Depth}
	{A low-risk vendor is assessed with a questionnaire. A high-risk vendor is also only assessed with a questionnaire. What is the issue?}
	{\item No issue
		\item Assessment depth should match vendor risk
		\item All vendors should be skipped
		\item All vendors need audits}
	{B}
	{\item Proportionality applies to vendor risk
		\item High-risk vendors need deeper review
		\item Uniform shallow assessment understates risk}
	{Assessment depth follows risk level.}
	
	\hq{89}{Subservice Organisations}
	{A SaaS vendor uses a subcontractor for storage. Which clause is MOST important?}
	{\item Subcontractor flow-down obligations and notification
		\item Vendor holiday party
		\item Vendor branding rules
		\item Pricing index}
	{A}
	{\item Flow-down carries obligations to subcontractors
		\item Notification provides visibility
		\item Other clauses do not address risk}
	{Flow-down obligations protect the chain.}
	
	\hq{90}{Vendor Exit}
	{The organization plans to end a contract. What is the FIRST step?}
	{\item Delete all data
		\item Confirm data return and destruction obligations plus transition plan
		\item Notify the press
		\item Cancel contracts immediately}
	{B}
	{\item Exit requires data obligations and transition
		\item Unsafe exit risks data loss or legal issues
		\item Immediate cancellation may breach terms}
	{Exit needs a plan, not a scramble.}
	
	% ============================================================
	\section{Governance and Risk Culture}
	% ============================================================
	
	\hq{91}{Risk Register Value}
	{Which outcome MOST shows a register is being used effectively?}
	{\item It grows in size each month
		\item It drives decisions, escalations and resource allocation
		\item It stays unchanged
		\item It has many pages}
	{B}
	{\item Use is proven by decisions
		\item Size and length do not prove value
		\item Static registers may be outdated}
	{Value is shown in decisions, not pages.}
	
	\hq{92}{Owner Motivation}
	{An owner refuses to accept the risk. What is the BEST practitioner response?}
	{\item Force acceptance
		\item Understand concerns, review the evidence and escalate if unresolved
		\item Ignore the owner
		\item Increase the rating}
	{B}
	{\item Owners may have context
		\item Review evidence addresses concerns
		\item Escalation resolves disagreements}
	{Understand and escalate, do not force.}
	
	\hq{93}{Culture Signals}
	{Which behavior MOST damages risk culture?}
	{\item Leaders openly discussing risks
		\item Leaders punishing those who escalate
		\item Encouraging early reporting
		\item Rewarding proactive risk management}
	{B}
	{\item Punishment discourages transparency
		\item Silence hides exposure
		\item Openness improves risk management}
	{Punishment kills honest reporting.}
	
	\hq{94}{Line Independence}
	{In the three lines model, which role should audit NOT perform?}
	{\item Independent assurance of controls
		\item Designing and operating controls
		\item Reporting to the board
		\item Following up on findings}
	{B}
	{\item Audit must remain independent
		\item Designing controls compromises independence
		\item Assurance is audit’s role}
	{Audit must not build what it audits.}
	
	\hq{95}{Practitioner Influence}
	{How does the practitioner BEST influence decisions?}
	{\item Make decisions for owners
		\item Provide clear analysis, options and consequences
		\item Withhold evidence
		\item Escalate every issue}
	{B}
	{\item Options and consequences support owners
		\item Transparency builds trust
		\item Withholding evidence damages credibility}
	{Influence through clarity, not control.}
	
	\hq{96}{New Tech Adoption}
	{The business wants to adopt a new AI tool urgently. What is the BEST governance response?}
	{\item Block immediately
		\item Conduct a proportionate risk assessment and approve or escalate
		\item Approve without review
		\item Delay indefinitely}
	{B}
	{\item Proportional assessment supports speed and safety
		\item Escalation follows risk thresholds
		\item Blocking and indefinite delay harm opportunity}
	{Proportionate assessment fits urgent adoption.}
	
	\hq{97}{Assessment Timing}
	{Which event should trigger a risk reassessment?}
	{\item A major change in leadership strategy
		\item A routine printer replacement
		\item A holiday announcement
		\item A marketing campaign}
	{A}
	{\item Strategic change alters risk
		\item Routine items rarely alter exposure
		\item Reassess when strategy changes}
	{Strategic change triggers reassessment.}
	
	\hq{98}{Executive Language}
	{A business leader says “we do not want to talk about IT issues.” What is the BEST practitioner response?}
	{\item Use technical jargon
		\item Reframe as business exposure, options and consequences
		\item Ignore the leader
		\item Escalate to the board}
	{B}
	{\item Business language enables engagement
		\item Reframing builds dialogue
		\item Escalation is for unresolved issues}
	{Speak business language to business leaders.}
	
	\hq{99}{Stale Data Risk}
	{A register entry was last reviewed three years ago. What is the BEST action?}
	{\item Assume still valid
		\item Review with the owner to confirm current status
		\item Delete the entry
		\item Increase the rating}
	{B}
	{\item Owners can confirm current state
		\item Review prevents misinformed decisions
		\item Assumption increases risk}
	{Confirm stale entries with owners.}
	
	\hq{100}{Owner Handover}
	{The owner of a critical risk is leaving the company. What should happen FIRST?}
	{\item Close the risk
		\item Hand over ownership formally to an appropriate successor
		\item Ignore the departure
		\item Lower the rating}
	{B}
	{\item Formal handover preserves accountability
		\item Successor must have authority
		\item Ignoring departure creates gaps}
	{Ownership must be handed over formally.}
	
	\begin{frame}{Exam Mindset Summary}
		\begin{block}{Remember}
			\begin{itemize}
				\item \textbf{FIRST:} understand context and assess before acting
				\item \textbf{Owners decide:} business owners accept risk; the practitioner advises
				\item \textbf{Appetite:} compare residual risk to appetite, then escalate or treat
				\item \textbf{Accountability} cannot be outsourced or transferred
				\item \textbf{Monitor:} validate data, trend it, and keep the register current
			\end{itemize}
		\end{block}
	\end{frame}
	
	\begin{frame}{Thank You}
		\begin{center}
			{\Huge \textbf{Good Luck on Your CRISC Exam!}}
		\end{center}
	\end{frame}
	
\end{document}